\documentclass[12pt,a4paper]{article}

\input{/home/tabaka/Documents/GitHub/Defs/defs.tex}

\setcounter{zadaniaRozdzial}{3}
\setcounter{section}{0}

\begin{document}
	
	\begin{center}
		\Large Wymagania edukacyjne z matematyki w klasie 2 Szkoły Branżowej na poszczególne oceny.
	\end{center}
	
	Wymagania edukacyjne na poszczególne oceny są stopniowane.
	Uzyskanie danej oceny wymaga spełnienia wymagań przewidzianych
	dla wszystkich ocen niższych, z uwzględnieniem wymagań
	określonych dla danej oceny.
	
	\section{Wyrażenia algebraiczne}
	
	\begin{longtable}{|p{16cm}|}
		\hline
		
		\textbf{Wymagania edukacyjne na ocenę dopuszczającą} \\ \hline
		
		\textbullet\ zna pojęcia jednomianu i sumy algebraicznej; \\
		
		\textbullet\ oblicza wartość liczbową wyrażeń algebraicznych w prostych przypadkach; \\
		
		\textbullet\ porządkuje jednomiany; \\
		
		\textbullet\ wypisuje wyrazy danej sumy algebraicznej; \\
		
		\textbullet\ wskazuje wyrazy podobne w danej sumie algebraicznej i przeprowadza ich redukcję; \\
		
		\textbullet\ dodaje i odejmuje sumy algebraiczne w prostych przypadkach; \\
		
		\textbullet\ mnoży sumy algebraiczne przez jednomiany; \\
		
		\textbullet\ wskazuje wspólny czynnik liczbowy wyrazów danej sumy algebraicznej; \\
		
		\textbullet\ zna wzory skróconego mnożenia: $(a+b)^2$, $(a-b)^2$, $a^2-b^2$. \\ \hline
		
		\textbf{Wymagania edukacyjne na ocenę dostateczną} \\ \hline
		
		\textbullet\ zapisuje w postaci wyrażeń algebraicznych zależności między zmiennymi opisanymi w zadaniach; \\
		
		\textbullet\ oblicza wartość liczbową wyrażeń algebraicznych; \\
		
		\textbullet\ zapisuje w postaci wyrażeń algebraicznych wyniki podanych działań; \\
		
		\textbullet\ dodaje i odejmuje sumy algebraiczne; \\
		
		\textbullet\ mnoży sumy algebraiczne składające się z dwóch składników i zapisuje wynik w najprostszej postaci; \\
		
		\textbullet\ stosuje odpowiedni wzór skróconego mnożenia do obliczania kwadratu sumy lub różnicy oraz różnicy kwadratów; \\
		
		\textbullet\ wyłącza jednomian przed nawias w sumie algebraicznej. \\ \hline
		
		\textbf{Wymagania edukacyjne na ocenę dobrą} \\ \hline
		
		\textbullet\ mnoży sumy algebraiczne i zapisuje wynik w najprostszej postaci; \\
		
		\textbullet\ zapisuje sumy algebraiczne w postaci iloczynu; \\
		
		\textbullet\ stosuje wyrażenia algebraiczne do obliczania obwodów i pól wielokątów; \\
		
		\textbullet\ stosuje wzory skróconego mnożenia do przekształcania wyrażeń algebraicznych. \\ \hline
		
		\textbf{Wymagania edukacyjne na ocenę bardzo dobrą} \\ \hline
		
		\textbullet\ stosuje działania na sumach algebraicznych w sytuacjach praktycznych; \\
		
		\textbullet\ stosuje metodę wyłączania wspólnego czynnika przed nawias do obliczania wartości wyrażeń arytmetycznych; \\
		
		\textbullet\ wykorzystuje wzory skróconego mnożenia do rozwiązywania problemów praktycznych. \\ \hline
		
		\textbf{Wymagania edukacyjne na ocenę celującą} \\ \hline
		
		\textbullet\ przeprowadza proste dowody, stosując działania na wyrażeniach algebraicznych oraz wzory skróconego mnożenia; \\
		
		\textbullet\ przygotowuje i prezentuje prace projektowe dotyczące zastosowania wyrażeń algebraicznych w sytuacjach praktycznych. \\ \hline
		
	\end{longtable}
	
	\section{Wykres funkcji kwadratowej}
	
	\begin{longtable}{|p{16cm}|}
		\hline
		
		\textbf{Wymagania edukacyjne na ocenę dopuszczającą} \\ \hline
		
		\textbullet\ zna pojęcia funkcji kwadratowej i paraboli; \\
		
		\textbullet\ szkicuje wykres funkcji $f(x)=ax^2$ i podaje jej własności; \\
		
		\textbullet\ interpretuje znak współczynnika $a$ funkcji kwadratowej; \\
		
		\textbullet\ szkicuje wykresy funkcji $f(x)=ax^2+q$ oraz $f(x)=a(x-p)^2$ i podaje ich własności; \\
		
		\textbullet\ sprawdza algebraicznie, czy dany punkt należy do wykresu funkcji kwadratowej; \\
		
		\textbullet\ rozpoznaje postać ogólną i kanoniczną funkcji kwadratowej; \\
		
		\textbullet\ odczytuje wartości współczynników funkcji kwadratowej zapisanej w postaci ogólnej; \\
		
		\textbullet\ oblicza wyróżnik funkcji kwadratowej; \\
		
		\textbullet\ przekształca wzór funkcji kwadratowej z postaci kanonicznej do postaci ogólnej; \\
		
		\textbullet\ odczytuje z wykresu funkcji kwadratowej jej najmniejszą lub największą wartość. \\ \hline
		
		\textbf{Wymagania edukacyjne na ocenę dostateczną} \\ \hline
		
		\textbullet\ szkicuje wykres funkcji kwadratowej zapisanej w postaci kanonicznej i podaje jej własności; \\
		
		\textbullet\ wyznacza wartość współczynnika $a$ funkcji $f(x)=ax^2$, gdy dane są współrzędne punktu należącego do jej wykresu; \\
		
		\textbullet\ ustala wzór funkcji kwadratowej na podstawie informacji o przesunięciach wykresu funkcji $f(x)=ax^2$ wzdłuż osi układu współrzędnych; \\
		
		\textbullet\ oblicza współrzędne wierzchołka paraboli i podaje równanie jej osi symetrii; \\
		
		\textbullet\ przekształca postać ogólną funkcji kwadratowej do postaci kanonicznej, stosując wzory na współrzędne wierzchołka; \\
		
		\textbullet\ oblicza najmniejszą lub największą wartość funkcji kwadratowej, zależnie od jej własności. \\ \hline
		
		\textbf{Wymagania edukacyjne na ocenę dobrą} \\ \hline
		
		\textbullet\ szkicuje wykres funkcji kwadratowej zapisanej w postaci ogólnej i podaje jej własności; \\
		
		\textbullet\ ustala wzór funkcji kwadratowej w postaci kanonicznej, gdy dane są współrzędne wierzchołka paraboli i innego punktu należącego do jej wykresu; \\
		
		\textbullet\ oblicza najmniejszą i największą wartość funkcji kwadratowej w przedziale domkniętym. \\ \hline
		
		\textbf{Wymagania edukacyjne na ocenę bardzo dobrą} \\ \hline
		
		\textbullet\ zapisuje wzór funkcji kwadratowej w postaci ogólnej na podstawie informacji o jej wykresie; \\
		
		\textbullet\ stosuje własności funkcji kwadratowej do rozwiązywania zadań optymalizacyjnych, w tym zadań osadzonych w kontekście praktycznym. \\ \hline
		
		\textbf{Wymagania edukacyjne na ocenę celującą} \\ \hline
		
		\textbullet\ rozwiązuje zadania o znacznym stopniu trudności dotyczące funkcji kwadratowej; \\
		
		\textbullet\ przygotowuje i prezentuje prace projektowe dotyczące zastosowania funkcji kwadratowej w sytuacjach praktycznych. \\ \hline
		
	\end{longtable}
	
	\section{Równanie i nierówności kwadratowe}
	
	\begin{longtable}{|p{16cm}|}
		\hline
		
		\textbf{Wymagania edukacyjne na ocenę dopuszczającą} \\ \hline
		
		\textbullet\ zna interpretację geometryczną rozwiązania równania kwadratowego; \\
		
		\textbullet\ odczytuje rozwiązania równania kwadratowego z wykresu odpowiedniej funkcji kwadratowej; \\
		
		\textbullet\ określa liczbę pierwiastków równania kwadratowego w zależności od znaku wyróżnika; \\
		
		\textbullet\ zna warunek istnienia postaci iloczynowej funkcji kwadratowej; \\
		
		\textbullet\ odczytuje miejsca zerowe funkcji kwadratowej z jej postaci iloczynowej; \\
		
		\textbullet\ zna związek między rozwiązaniem nierówności kwadratowej a znakiem wartości odpowiedniej funkcji kwadratowej. \\ \hline
		
		\textbf{Wymagania edukacyjne na ocenę dostateczną} \\ \hline
		
		\textbullet\ rozwiązuje proste równania kwadratowe metodą rozkładu na czynniki oraz z zastosowaniem wzorów skróconego mnożenia; \\
		
		\textbullet\ rozwiązuje równania kwadratowe, stosując wzory na pierwiastki; \\
		
		\textbullet\ przedstawia funkcję kwadratową w postaci iloczynowej, jeśli taka postać istnieje; \\
		
		\textbullet\ stosuje zależność między miejscami zerowymi funkcji kwadratowej a pierwszą współrzędną wierzchołka paraboli; \\
		
		\textbullet\ odczytuje z wykresu funkcji kwadratowej $f$ zbiór rozwiązań nierówności typu $f(x)>0$, $f(x)<0$, $f(x)\leq 0$ lub $f(x)\geq 0$. \\ \hline
		
		\textbf{Wymagania edukacyjne na ocenę dobrą} \\ \hline
		
		\textbullet\ wyznacza algebraicznie współrzędne punktów przecięcia paraboli z osiami układu współrzędnych i na tej podstawie szkicuje jej wykres; \\
		
		\textbullet\ rozwiązuje nierówności kwadratowe; \\
		
		\textbullet\ wyznacza współczynniki funkcji kwadratowej na podstawie informacji o jej własnościach w prostych przypadkach. \\ \hline
		
		\textbf{Wymagania edukacyjne na ocenę bardzo dobrą} \\ \hline
		
		\textbullet\ stosuje równania kwadratowe do rozwiązywania zadań tekstowych; \\
		
		\textbullet\ analizuje zadanie tekstowe, zapisuje odpowiednią nierówność kwadratową opisującą daną zależność i rozwiązuje ją; \\
		
		\textbullet\ wyznacza współczynniki funkcji kwadratowej na podstawie informacji o jej własnościach. \\ \hline
		
		\textbf{Wymagania edukacyjne na ocenę celującą} \\ \hline
		
		\textbullet\ rozwiązuje zadania o podwyższonym stopniu trudności. \\ \hline
		
	\end{longtable}
	
	\section{Wielokąty}
	
	\begin{longtable}{|p{16cm}|}
		\hline
		
		\textbf{Wymagania edukacyjne na ocenę dopuszczającą} \\ \hline
		
		\textbullet\ zna klasyfikację trójkątów ze względu na miary kątów oraz długości boków; \\
		
		\textbullet\ stosuje twierdzenie o sumie miar kątów wewnętrznych trójkąta; \\
		
		\textbullet\ oblicza miary kątów trójkąta, stosując własności kątów wierzchołkowych i przyległych; \\
		
		\textbullet\ zna pojęcia środkowej trójkąta, ortocentrum i środka ciężkości trójkąta oraz wskazuje je w danym trójkącie; \\
		
		\textbullet\ zna cechy przystawania trójkątów; \\
		
		\textbullet\ sprawdza, czy z odcinków o podanych długościach można zbudować trójkąt, stosując warunek istnienia trójkąta; \\
		
		\textbullet\ zna twierdzenie Pitagorasa i twierdzenie do niego odwrotne; \\
		
		\textbullet\ stosuje wzór na pole trójkąta; \\
		
		\textbullet\ zna własności kątów i przekątnych w kwadratach, prostokątach, równoległobokach, rombach i trapezach; \\
		
		\textbullet\ stosuje wzory na pole kwadratu, prostokąta, równoległoboku, rombu i trapezu. \\ \hline
		
		\textbf{Wymagania edukacyjne na ocenę dostateczną} \\ \hline
		
		\textbullet\ oblicza miary kątów trójkąta, stosując własności kątów odpowiadających i naprzemianległych; \\
		
		\textbullet\ konstruuje wysokość i środkową trójkąta; \\
		
		\textbullet\ rozpoznaje trójkąty przystające; \\
		
		\textbullet\ uzasadnia przystawanie trójkątów, powołując się na odpowiednią cechę; \\
		
		\textbullet\ stosuje twierdzenie Pitagorasa do wyznaczania długości boków trójkąta prostokątnego; \\
		
		\textbullet\ stosuje twierdzenie odwrotne do twierdzenia Pitagorasa w prostych przypadkach; \\
		
		\textbullet\ oblicza odległość punktów w układzie współrzędnych; \\
		
		\textbullet\ korzysta ze wzorów na długość przekątnej kwadratu i wysokość trójkąta równobocznego; \\
		
		\textbullet\ wyznacza brakujące długości boków w trójkątach o kątach $45^\circ$, $45^\circ$, $90^\circ$ oraz $30^\circ$, $60^\circ$, $90^\circ$, gdy zna długość jednego z boków; \\
		
		\textbullet\ stosuje wzór na pole trójkąta równobocznego $P=\frac{a^2\sqrt{3}}{4}$; \\
		
		\textbullet\ stosuje własności kątów i przekątnych w kwadratach, prostokątach, równoległobokach, rombach i trapezach do obliczania ich obwodów i pól. \\ \hline
		
		\textbf{Wymagania edukacyjne na ocenę dobrą} \\ \hline
		
		\textbullet\ stosuje własności wysokości trójkąta do rozwiązywania zadań; \\
		
		\textbullet\ stosuje własności środkowych trójkąta do rozwiązywania zadań; \\
		
		\textbullet\ wyznacza konstrukcyjnie środek ciężkości trójkąta; \\
		
		\textbullet\ stosuje twierdzenie Pitagorasa do wyznaczania długości odcinków w wielokątach; \\
		
		\textbullet\ stosuje twierdzenie odwrotne do twierdzenia Pitagorasa; \\
		
		\textbullet\ oblicza obwód wielokąta, mając dane współrzędne jego wierzchołków. \\ \hline
		
		\textbf{Wymagania edukacyjne na ocenę bardzo dobrą} \\ \hline
		
		\textbullet\ stosuje własności wielokątów do rozwiązywania zadań osadzonych w kontekście praktycznym; \\
		
		\textbullet\ wyznacza konstrukcyjnie ortocentrum trójkąta; \\
		
		\textbullet\ stosuje własności ortocentrum i środka ciężkości trójkąta w sytuacjach praktycznych. \\ \hline
		
		\textbf{Wymagania edukacyjne na ocenę celującą} \\ \hline
		
		\textbullet\ rozwiązuje zadania o podwyższonym stopniu trudności; \\
		
		\textbullet\ przygotowuje i prezentuje prace projektowe dotyczące zastosowania wielokątów w sytuacjach praktycznych. \\ \hline
		
	\end{longtable}
	
	\section{Podobieństwo}
	
	\begin{longtable}{|p{16cm}|}
		\hline
		
		\textbf{Wymagania edukacyjne na ocenę dopuszczającą} \\ \hline
		
		\textbullet\ rozpoznaje figury podobne; \\
		
		\textbullet\ podaje cechy podobieństwa trójkątów; \\
		
		\textbullet\ zna zależność między długościami boków wielokątów podobnych a skalą podobieństwa; \\
		
		\textbullet\ zna zależność między polami wielokątów podobnych a skalą podobieństwa. \\ \hline
		
		\textbf{Wymagania edukacyjne na ocenę dostateczną} \\ \hline
		
		\textbullet\ uzasadnia podobieństwo trójkątów, powołując się na odpowiednią cechę podobieństwa; \\
		
		\textbullet\ układa odpowiednią proporcję w celu wyznaczenia szukanych długości boków trójkątów podobnych; \\
		
		\textbullet\ oblicza długości boków wielokątów podobnych w danej skali; \\
		
		\textbullet\ ustala skalę podobieństwa wielokątów podobnych, gdy dane są długości ich boków lub pola. \\ \hline
		
		\textbf{Wymagania edukacyjne na ocenę dobrą} \\ \hline
		
		\textbullet\ stosuje zależności między obwodami wielokątów podobnych a skalą podobieństwa do rozwiązywania zadań; \\
		
		\textbullet\ stosuje zależności między polami wielokątów podobnych a skalą podobieństwa do rozwiązywania zadań. \\ \hline
		
		\textbf{Wymagania edukacyjne na ocenę bardzo dobrą} \\ \hline
		
		\textbullet\ stosuje własności wielokątów podobnych do rozwiązywania zadań osadzonych w kontekście praktycznym. \\ \hline
		
		\textbf{Wymagania edukacyjne na ocenę celującą} \\ \hline
		
		\textbullet\ rozwiązuje zadania o podwyższonym stopniu trudności; \\
		
		\textbullet\ przygotowuje i prezentuje prace projektowe dotyczące zastosowania podobieństwa figur w sytuacjach praktycznych, np. wykonuje modele i makiety oraz przeprowadza pomiary wysokości niedostępnych obiektów. \\ \hline
		
	\end{longtable}
	
	\section{Trygonometria}
	
	\begin{longtable}{|p{16cm}|}
		\hline
		
		\textbf{Wymagania edukacyjne na ocenę dopuszczającą} \\ \hline
		
		\textbullet\ zna definicje funkcji trygonometrycznych kąta ostrego w trójkącie prostokątnym; \\
		
		\textbullet\ podaje wartości funkcji trygonometrycznych kątów $30^\circ$, $45^\circ$ i $60^\circ$; \\
		
		\textbullet\ odczytuje wartości funkcji trygonometrycznych danego kąta ostrego, korzystając z tablic lub kalkulatora; \\
		
		\textbullet\ zna związki między funkcjami trygonometrycznymi tego samego kąta: $\sin^2\alpha+\cos^2\alpha=1$ oraz $\tg\alpha=\frac{\sin\alpha}{\cos\alpha}$. \\ \hline
		
		\textbf{Wymagania edukacyjne na ocenę dostateczną} \\ \hline
		
		\textbullet\ wyznacza wartości funkcji trygonometrycznych kątów $30^\circ$, $45^\circ$ i $60^\circ$; \\
		
		\textbullet\ oblicza wartości funkcji trygonometrycznych w trójkącie prostokątnym o podanych długościach boków; \\
		
		\textbullet\ odczytuje z tablic lub wyznacza za pomocą kalkulatora miarę kąta ostrego, gdy zna wartość jego funkcji trygonometrycznej; \\
		
		\textbullet\ oblicza miary kątów w trójkącie prostokątnym o podanych długościach boków; \\
		
		\textbullet\ wyznacza długości boków trójkąta prostokątnego za pomocą funkcji trygonometrycznych; \\
		
		\textbullet\ stosuje funkcje trygonometryczne do wyznaczania długości odcinków w figurach płaskich oraz obliczania ich pól. \\ \hline
		
		\textbf{Wymagania edukacyjne na ocenę dobrą} \\ \hline
		
		\textbullet\ oblicza wartości pozostałych funkcji trygonometrycznych kąta ostrego, gdy dana jest wartość jego sinusa lub cosinusa; \\
		
		\textbullet\ stosuje wzór na pole trójkąta $P=\frac{1}{2}ab\sin\alpha$ oraz wzór na pole równoległoboku $P=ab\sin\alpha$ w zadaniach; \\
		
		\textbullet\ korzysta ze wzorów redukcyjnych w zadaniach; \\
		
		\textbullet\ stosuje poznane związki między funkcjami trygonometrycznymi do upraszczania wyrażeń zawierających te funkcje w prostych przypadkach. \\ \hline
		
		\textbf{Wymagania edukacyjne na ocenę bardzo dobrą} \\ \hline
		
		\textbullet\ oblicza wartości pozostałych funkcji trygonometrycznych kąta ostrego, gdy dana jest wartość jego tangensa; \\
		
		\textbullet\ oblicza pole wielokąta, wykorzystując wzór na pole trójkąta; \\
		
		\textbullet\ stosuje funkcje trygonometryczne do rozwiązywania zadań praktycznych, w tym wykorzystuje zależność między tangensem kąta a nachyleniem drogi; \\
		
		\textbullet\ stosuje poznane związki między funkcjami trygonometrycznymi do upraszczania wyrażeń zawierających te funkcje. \\ \hline
		
		\textbf{Wymagania edukacyjne na ocenę celującą} \\ \hline
		
		\textbullet\ rozwiązuje zadania o podwyższonym stopniu trudności; \\
		
		\textbullet\ uzasadnia związki między funkcjami trygonometrycznymi tego samego kąta oraz zależności występujące we wzorach redukcyjnych; \\
		
		\textbullet\ przygotowuje i prezentuje prace projektowe dotyczące zastosowania funkcji trygonometrycznych w sytuacjach praktycznych. \\ \hline
		
	\end{longtable}
	
	\section{Okręgi i wielokąty}
	
	\begin{longtable}{|p{16cm}|}
		\hline
		
		\textbf{Wymagania edukacyjne na ocenę dopuszczającą} \\ \hline
		
		\textbullet\ zna i stosuje wzory na długość okręgu i pole koła; \\
		
		\textbullet\ rozpoznaje kąty środkowe w okręgu oraz wskazuje łuki, na których są one oparte; \\
		
		\textbullet\ rozpoznaje kąty wpisane w okrąg oraz wskazuje łuki, na których są one oparte; \\
		
		\textbullet\ zna twierdzenie o kącie środkowym i kącie wpisanym opartych na tym samym łuku oraz wnioski z tego twierdzenia; \\
		
		\textbullet\ wyznacza liczbę punktów wspólnych prostej i okręgu; \\
		
		\textbullet\ wskazuje w danym trójkącie środek okręgu wpisanego w ten trójkąt i środek okręgu opisanego na nim oraz podaje własności tych punktów; \\
		
		\textbullet\ rozpoznaje wielokąty foremne i zna ich własności; \\
		
		\textbullet\ ustala, czy dany wielokąt foremny ma środek symetrii; \\
		
		\textbullet\ podaje liczbę osi symetrii wielokąta foremnego. \\ \hline
		
		\textbf{Wymagania edukacyjne na ocenę dostateczną} \\ \hline
		
		\textbullet\ wyznacza promień okręgu o danej długości; \\
		
		\textbullet\ wyznacza promień koła o danym polu; \\
		
		\textbullet\ oblicza pole pierścienia kołowego; \\
		
		\textbullet\ oblicza długość łuku i pole wycinka koła wyznaczonego przez dany kąt środkowy; \\
		
		\textbullet\ wyznacza miary kątów środkowych i wpisanych opartych na tym samym łuku; \\
		
		\textbullet\ rozpoznaje kąty wpisane oparte na tym samym łuku i korzysta z równości ich miar; \\
		
		\textbullet\ korzysta z własności kąta wpisanego opartego na półokręgu w prostych przypadkach; \\
		
		\textbullet\ określa wzajemne położenie okręgu i prostej, porównując odległość środka okręgu od prostej z długością promienia okręgu; \\
		
		\textbullet\ ustala położenie środka okręgu opisanego na trójkącie na podstawie informacji o jego kątach; \\
		
		\textbullet\ konstruuje symetralną odcinka i dwusieczną kąta; \\
		
		\textbullet\ oblicza długość promienia okręgu wpisanego w kwadrat lub okręgu opisanego na kwadracie o danej długości boku. \\ \hline
		
		\textbf{Wymagania edukacyjne na ocenę dobrą} \\ \hline
		
		\textbullet\ stosuje wzory na długość okręgu i pole koła w sytuacjach praktycznych; \\
		
		\textbullet\ oblicza miarę kąta środkowego, gdy dana jest długość łuku lub pole wycinka koła wyznaczonego przez ten kąt; \\
		
		\textbullet\ korzysta z własności kąta wpisanego opartego na półokręgu; \\
		
		\textbullet\ konstruuje styczną do okręgu; \\
		
		\textbullet\ korzysta z własności stycznej do okręgu; \\
		
		\textbullet\ stosuje własności środka okręgu wpisanego w trójkąt; \\
		
		\textbullet\ stosuje własności środka okręgu opisanego na trójkącie; \\
		
		\textbullet\ oblicza długość promienia okręgu wpisanego w trójkąt równoboczny lub sześciokąt foremny o danej długości boku; \\
		
		\textbullet\ oblicza długość promienia okręgu opisanego na trójkącie równobocznym lub sześciokącie foremnym o danej długości boku; \\
		
		\textbullet\ wyznacza długość boku kwadratu, gdy dana jest długość promienia okręgu wpisanego w kwadrat lub okręgu opisanego na kwadracie; \\
		
		\textbullet\ wyznacza miarę kąta wewnętrznego wielokąta foremnego. \\ \hline
		
		\textbf{Wymagania edukacyjne na ocenę bardzo dobrą} \\ \hline
		
		\textbullet\ wyznacza pola odcinków kołowych; \\
		
		\textbullet\ korzysta z twierdzenia o odcinkach stycznych; \\
		
		\textbullet\ konstruuje okrąg wpisany w trójkąt lub okrąg opisany na trójkącie; \\
		
		\textbullet\ oblicza długość boku trójkąta równobocznego lub sześciokąta foremnego, gdy dana jest długość promienia okręgu wpisanego w ten wielokąt lub okręgu opisanego na nim; \\
		
		\textbullet\ stosuje zależności między długością boku trójkąta równobocznego, kwadratu lub sześciokąta foremnego a promieniem okręgu opisanego na tych wielokątach lub okręgu w nie wpisanego w sytuacjach praktycznych; \\
		
		\textbullet\ wyznacza liczbę boków wielokąta foremnego, gdy dana jest miara jego kąta wewnętrznego. \\ \hline
		
		\textbf{Wymagania edukacyjne na ocenę celującą} \\ \hline
		
		\textbullet\ udowadnia zależności między długością boku kwadratu lub trójkąta równobocznego a długością promienia okręgu opisanego na tych wielokątach lub okręgu w nie wpisanego; \\
		
		\textbullet\ przygotowuje i prezentuje prace projektowe dotyczące zastosowania kół, okręgów lub wielokątów w sytuacjach praktycznych. \\ \hline
		
	\end{longtable}
	
\end{document}